% !TEX root = ../print.tex
% Draft \S2 --- Preliminaries and notation
%
% Source: draft/hemigroup-causal-scale-space-kernels.md, lines 31-67.
% Chapter 2 of the blueprint; the shared counter puts every number here on the draft's own.
%
% TWO DEPARTURES FROM THE DRAFT, both deliberate:
%
% 1. Proposition 2.3 stays ONE node with four parts, rather than four nodes. Splitting it
%    would push Lemma 2.4 to 2.7 and break the numbering alignment for the section the rest
%    of the article cites most. The per-clause honesty lives in the \textbf{Assignment.}
%    clause instead, which is what that mechanism is for (LINKAGE.md rule 7).
%
% 2. prop:laplace-continuity (2.5) and prop:laplace-uniqueness (2.6) are ADDITIVE -- the
%    draft uses both facts (Lemma 4.1's uniqueness clause, Lemma 6.1, Theorem 7.3's (A7)
%    verification, Proposition 9.2, Theorem 9.5, Lemma 10.3(5)) but never states them, so
%    their trust was invisible to the dependency graph. Appending them leaves 2.1-2.4 alone.
%
% 3. 2.7-2.10 are ADDITIVE for the same reason and by the same rule -- appended, so that
%    2.1-2.6 keep the draft's numbers. They exist because the FORMALISATION had outrun the
%    graph: ~900 lines of finished, axiom-guarded Lean (Levy, Injectivity, WeakConvergence,
%    the Markov bound of Continuity) proved things this chapter either did not state at all
%    (the representation, Def. 2.7 -- the single most-used object in the Lean development,
%    and until now nameable only inside Prop. 2.3(3)) or stated only in a stronger cited
%    form it could not be tagged against (2.5, 2.6). The pattern for the latter is worth
%    naming, because chapters 8-12 will want it again:
%
%      an [A] node keeps the PAPER's statement and its ledger entry; a [T] node beneath it
%      states the RESTRICTED case the development actually proves, carries the \lean tag,
%      and is what the downstream node \uses.
%
%    So Prop. 2.5 and 2.6 are unchanged -- the paper still cites Feller -- while 2.8 and 2.9
%    record that nothing formalised here rests on A5 or A6.
% \notes{} tags are omitted throughout until the hub's content-note slugs exist; an
% unresolvable slug fails `linkage check --wiki`.

Time runs over $\RR$. We write $X = \Lone$, real, with positive cone $X_+$ and
$\int f := \int_{\RR} f(t)\,dt$; $\Mplus(\RR)$ for the finite positive Borel measures on
$\RR$, with total mass $\|\mu\| = \mu(\RR)$; translations $(\shift_a f)(t) = f(t-a)$ for
$a \in \RR$; dilations on functions $(\dil_\sigma f)(t) = \sigma^{-1} f(t/\sigma)$,
$\sigma > 0$, which preserve mass, and on measures $\dil_\sigma \mu :=$ the pushforward of
$\mu$ under $t \mapsto \sigma t$; and convolution
$(\mu * f)(t) = \int f(t - r)\,\mu(dr)$ for $\mu \in \Mplus(\RR)$, $f \in X$. Then
$\|\mu * f\|_1 \le \|\mu\|\,\|f\|_1$, and dilation is compatible with convolution:
\[
  \dil_\sigma(\mu * f) = (\dil_\sigma \mu) * (\dil_\sigma f).
\]

For $\mu \in \Mplus$ supported on $[0,\infty)$ the Laplace transform is
\[
  \hat\mu(s) = \int_{[0,\infty)} e^{-st}\,\mu(dt), \qquad s \ge 0 .
\]
If $\mu \ne 0$ then $\hat\mu(s) > 0$ for all $s \ge 0$.

% draft: Definition 2.1
\begin{definition}[completely monotone]\label{def:completely-monotone}
A function $f : (0,\infty) \to \RR$ is \emph{completely monotone} (CM) if $f \in C^\infty$
and $(-1)^n f^{(n)} \ge 0$ for all $n \ge 0$.
\statusT
\end{definition}

% draft: Definition 2.2
\begin{definition}[Bernstein function]\label{def:bernstein-function}
\uses{def:completely-monotone}
A function $g : (0,\infty) \to [0,\infty)$ is a \emph{Bernstein function} if $g \in C^\infty$
and $g'$ is completely monotone. We write $g \in \BF$, and $g \in \BFz$ if moreover
$g(\zp) = 0$.
\statusT
\end{definition}

% draft: Proposition 2.3 (toolbox)
\begin{proposition}[the Bernstein-function toolbox]\label{prop:bernstein-toolbox}
\uses{def:completely-monotone,def:bernstein-function}
\begin{enumerate}
\item (Bernstein--Widder / Feller criterion.) $f$ is the Laplace transform of a probability
measure on $[0,\infty)$ if and only if $f$ is completely monotone and $f(\zp) = 1$. More
generally, $f$ is completely monotone if and only if $f(s) = \int_{[0,\infty)} e^{-st}\,F(dt)$
for a measure $F$ on $[0,\infty)$, not necessarily finite, and $F$ is then unique.
\item For $g : (0,\infty) \to [0,\infty)$ the following are equivalent: $g \in \BF$; the
composition $h \circ g$ is completely monotone for every completely monotone $h$; and
$e^{-\tau g}$ is completely monotone for every $\tau > 0$.
\item Every $g \in \BF$ is nonnegative, nondecreasing and concave, and has the
Lévy--Khintchine representation
\[
  g(s) = a + b s + \int_0^\infty \bigl(1 - e^{-st}\bigr)\,\nu(dt), \qquad
  a, b \ge 0, \quad \int_0^\infty (1 \wedge t)\,\nu(dt) < \infty,
\]
with $g(\zp) = a$. The triple $(a, b, \nu)$ is unique.
\item $\BF$ is a convex cone, closed under pointwise limits: if $g_n \in \BF$ and
$g_n(s) \to g(s) < \infty$ for every $s > 0$, then $g \in \BF$. Consequently $\BF$ is closed
under mixtures $s \mapsto \int g_u(s)\,m(du)$ with $m \ge 0$, whenever the integral is finite.
\end{enumerate}
\statusA\quad\ledger{A1}\quad\ledger{A2}\quad\ledger{A3}\quad\ledger{A4}\quad\emph{(the
classical Bernstein-function toolbox: \sscite{feller2009introduction}, Ch.~XIII and
\sscite{schilling2012bernstein}, Ch.~1 and~3. Cited, not reproven here.
\textbf{Assignment.} The four parts are carried by four separate entries, one each.
Part~(1) is \ledger{A1} entire, in both the probability form and the general measure form
with uniqueness. Part~(2) is \ledger{A2} entire --- all three equivalences of the cited
theorem, not only the exponential one: the composition clause is what
\ref{lem:potential-kernel} inverts a Bernstein function with. Part~(3) is \ledger{A3} entire, including
the uniqueness of the triple, which Lemma~7.1 needs and which a representation theorem
without it would not supply. Part~(3) is also the \emph{bridge} between the two vocabularies
this project speaks: it is what identifies \ref{def:bernstein-function} with
\ref{def:levy-exponent}, and therefore the single edge every Lean tag in chapters 5--7 crosses,
since the development names admissibility by the representation and never by derivative signs.
A node whose conclusion is stated in $\BFz$ cannot carry a Lean tag without spending this part;
a node whose conclusion is stated in $\LE$ carries one for free.
Part~(4) is \ledger{A4} entire; it is the single external fact
Theorem~5.2 rests on, and the null-array argument that produces the pointwise limit there is
ours and is held \textbf{[T]}. None of the four carries anything about the hemigroup family:
they are statements about $\CM$ and $\BF$ as classes of functions, and every step joining
them to $\Phi_{x,y}$ is \textbf{[T]}.)}
\end{proposition}

% draft: Lemma 2.4
\begin{lemma}[vanishing lemma]\label{lem:vanishing}
\uses{def:bernstein-function,prop:bernstein-toolbox}
\lean{Hemigroup.levyExponent_eq_zero_of_eq_zero}\leanok
If $g \in \BFz$ and $g(s_0) = 0$ for some $s_0 > 0$, then $g \equiv 0$.
\statusT\quad\emph{(the Lean development states this for the \emph{representation}
$g(s) = b_0 s + \int (1 - e^{-st})\,\nu(dt)$ rather than for the derivative-sign definition
\ref{def:bernstein-function}, following the formalisation strategy; the two notions agree by
the Lévy--Khintchine theorem, which is A3 and reaches this node through
\ref{prop:bernstein-toolbox}(3). From the representation the proof is shorter than the one
below --- a sum of nonnegative terms vanishes only when each does --- and needs neither
concavity nor complete monotonicity. \texttt{\#print axioms} reduces to Lean core only.)}
\end{lemma}

\begin{proof}
\leanok
$g$ is nonnegative and concave by \ref{prop:bernstein-toolbox}(3), with
$g(\zp) = g(s_0) = 0$. A nonnegative concave function vanishing at both ends of $[0, s_0]$
vanishes on it, so $g \equiv 0$ on $(0, s_0]$ and hence $g' \equiv 0$ on $(0, s_0)$. Now $g'$
is nonincreasing, being the derivative of a concave function, and nonnegative, being
completely monotone; a nonnegative nonincreasing function that vanishes on $(0,s_0)$ vanishes
identically. So $g' \equiv 0$ on $(0,\infty)$, and with $g(\zp) = 0$ this gives $g \equiv 0$.
\end{proof}

\bigskip

The next two statements are not in the draft, which uses both facts without stating either.
They are recorded here so that the trust they carry is visible in the dependency graph rather
than buried inside proofs. Both are classical, and both are cited to the same section of
\sscite{feller2009introduction}; note that that section states the probability and the measure
forms as four separate theorems, on four different pages, so the pair a given proof needs must
be named individually.

% additive: not in the draft. Used by lem:convolution-representation, thm:main-characterization.
\begin{proposition}[continuity theorem for Laplace transforms]\label{prop:laplace-continuity}
Let $\mu_n, \mu$ be probability measures on $[0,\infty)$. Then $\mu_n \to \mu$ weakly if and
only if $\hat\mu_n(s) \to \hat\mu(s)$ for every $s > 0$. The same equivalence holds for
arbitrary measures on $[0,\infty)$, provided that in the direction from transforms to measures
the masses $\hat\mu_n(a)$ are bounded for some $a > 0$.
\statusA\quad\ledger{A5}\quad\emph{(Feller's continuity theorem;
\sscite{feller2009introduction}, Vol.~2, \S XIII.1, Theorem~2 and Theorem~2a. A cited
interface, not reproven here. \textbf{Assignment.} \ledger{A5} carries both clauses. The
boundedness hypothesis in the second is not decoration: Feller gives a counterexample on the
same page showing it cannot be dropped, and any proof here that invokes the measure form must
discharge it explicitly. The probability form carries it for free, which is why Theorem~7.3
uses that one. What \ledger{A5} does not carry is the tightness argument that supplies weak
convergence in the first place; that is \textbf{[T]}. \textbf{Not needed in Lean} (resolved
2026-08-09): the development proves the form it uses rather than citing it, as
\texttt{Hemigroup.tendsto\_integral\_of\_tendsto\_laplace}, whose \texttt{\#print axioms} reduces
to Lean core. For \emph{causal probability} measures, convergence of the transforms together with
tightness gives weak convergence; that is precisely what Theorem~7.3$(\Leftarrow)$ uses, and the
boundedness hypothesis flagged above is carried for free by the probability form, so the trap is
designed out rather than remembered. The general measure form is deliberately \emph{not}
formalised. Tightness is ours, as this clause always claimed and now checks: the Markov bound
$\mu(t>T)(1-e^{-sT}) \le 1 - \hat\mu(s)$ (\texttt{measureReal\_Ioi\_mul\_le}), uniform over the
family as $\mu_{a,b}(t>T) \le F(Bs)/(1-e^{-sT})$ (\texttt{kernel\_tail\_le}), with $F(\zp) = 0$
(\texttt{exists\_exponent\_lt}) making $s$ choosable. The node keeps \textbf{[A]} because that is
the \emph{paper's} stance, and carries no Lean tag because what is proved is one direction of the
probability case, not the statement above.)}
\end{proposition}

% additive: not in the draft. Used by lem:convolution-representation, lem:covariance-laplace,
% prop:scale-evolution, thm:sonine-conservation, lem:generator-properties.
\begin{proposition}[uniqueness of the Laplace transform]\label{prop:laplace-uniqueness}
A measure on $[0,\infty)$ is determined by its Laplace transform: if $\hat\mu(s) = \hat\rho(s)$
for all $s$ in some interval $(a,\infty)$, then $\mu = \rho$.
\statusA\quad\ledger{A6}\quad\emph{(Feller's uniqueness theorem;
\sscite{feller2009introduction}, Vol.~2, \S XIII.1, Theorem~1 and Theorem~1a. A cited
interface, not reproven here. \textbf{Assignment.} \ledger{A6} carries the statement entire,
for probability measures (Theorem~1) and for arbitrary measures (Theorem~1a). This is what
licenses every step in the article of the form ``the transforms agree, therefore the measures
agree'' --- the uniqueness clause of Lemma~4.1, and Lemma~6.1, Proposition~9.2, Theorem~9.5
and Lemma~10.3(5). It says nothing about existence, and nothing about which functions are
transforms; that is \ref{prop:bernstein-toolbox}(1). \textbf{Not needed in Lean.} The
development proves this rather than citing it: \texttt{Hemigroup.laplace\_injective}, by the
classical substitution $x = e^{-t}$, which carries a causal measure onto $[0,1]$, turns
$\hat\mu(n)$ into the $n$th moment of the image, and closes with the Weierstrass approximation
theorem. \texttt{\#print axioms} reduces to Lean core, so A6 never enters
\texttt{trust-boundary.txt} and nothing proved in Lean rests on it. The node keeps \textbf{[A]}
because that is the \emph{paper's} stance, and it carries no Lean tag because what is proved is
weaker than what is stated: agreement on all of $[0,\infty)$, not on a tail $(a,\infty)$. That
gap is not a formality --- equality of the moments from some $n_0$ onwards does not by itself
pin down an atom at the origin of the transported measure, and the argument that ours has none
is not formalised.)}
\end{proposition}

\bigskip

The remaining four statements are additive too, and they are the \emph{formalisation's}
vocabulary rather than the draft's. \ref{def:levy-exponent} names the object the Lean
development is written in; the three after it are the restricted cases of \ref{prop:laplace-continuity}
and \ref{prop:laplace-uniqueness} that it proves outright, recorded as \textbf{[T]} nodes beneath
the \textbf{[A]} ones they refine so that the graph shows what actually rests on \ledger{A5} and
\ledger{A6}. Nothing formalised does.

% additive: not in the draft. The representation as a named class -- what the Lean
% development speaks. Used by lem:vanishing, thm:increments-bernstein, prop:canonical-gauge,
% lem:selfdecomposable-exponents, thm:main-construction.
\begin{definition}[Lévy exponent]\label{def:levy-exponent}
\uses{prop:bernstein-toolbox}
\lean{Hemigroup.levyExponent}\leanok
A function $g : [0,\infty) \to [0,\infty]$ is a \emph{Lévy exponent} if
\[
  g(s) = b_0\,s + \int_0^\infty \bigl(1 - e^{-st}\bigr)\,\nu(dt),
  \qquad b_0 \ge 0,
\]
for a drift $b_0$ and a positive Borel measure $\nu$ on $(0,\infty)$ --- with no killing term,
so that $g(0) = 0$. We write $g \in \LE$, and call $(b_0, \nu)$ its \emph{triple}.
\statusT\quad\emph{(A \emph{name} for \ref{prop:bernstein-toolbox}(3)'s right-hand side, not a
new fact: by that part, $\LE = \BFz$ as classes of functions, with the triple unique. Both names
are kept because the two documents need different ones. The paper argues in $\BFz$, where the
draft's proofs live; the Lean development argues in $\LE$ and never defines
\ref{def:completely-monotone} at all --- the reason is \texttt{blueprint/DESIGN-formalization-strategy.md},
and the payoff is that the Lévy--Khintchine theorem (\ledger{A3}) is crossed once, here, instead
of inside every proof. Stated in Lean as \texttt{Hemigroup.levyExponent}, valued in
$[0,\infty]$ so that no integrability hypothesis is carried; the integrability condition
$\int (1 \wedge t)\,\nu(dt) < \infty$ of \ref{prop:bernstein-toolbox}(3) is equivalent, for
causal $\nu$, to finiteness of $g$ at a single $s > 0$, and appears in the development as that
hypothesis instead. \ref{lem:vanishing} and \ref{lem:selfdecomposable-increment}
are both proved directly from this definition
(\texttt{levyExponent\_eq\_zero\_of\_eq\_zero}, \texttt{levyExponentD\_increment}).)}
\end{definition}

% additive: not in the draft. The [T] refinement of prop:laplace-uniqueness.
\begin{proposition}[uniqueness of the Laplace transform, causal case]\label{prop:laplace-uniqueness-causal}
\uses{prop:laplace-uniqueness}
\lean{Hemigroup.laplace_injective}\leanok
Let $\mu, \rho$ be finite measures on $\RR$ carried by $[0,\infty)$. If
$\hat\mu(s) = \hat\rho(s)$ for every $s \ge 0$, then $\mu = \rho$.
\statusT\quad\emph{(the case of \ref{prop:laplace-uniqueness} that this article's own arguments
use, and it is \emph{proved}: \texttt{Hemigroup.laplace\_injective},
\texttt{\#print axioms} reducing to Lean core. It is weaker than the cited statement in exactly
one respect --- agreement is assumed on all of $[0,\infty)$, not on a tail $(a,\infty)$ --- and
that gap is not a formality, which is why \ref{prop:laplace-uniqueness} keeps its \ledger{A6}
tag and this node does not replace it. Every step in the formalised part of the article of the
form ``the transforms agree, therefore the measures agree'' has both measures carried by
$[0,\infty)$ and both transforms known on all of it, so it cites this node and not that one; the
consequence is that \ledger{A6} does not appear in \texttt{blueprint/trust-boundary.txt}.)}

\emph{\textbf{Extended 2026-08-11 to measures that are not finite}, which is what Chapter~9
needs: $\kappa^{(x)}$, $\ell^{(x)}$ and Lebesgue measure are none of them finite, and
\ref{thm:sonine-conservation} compares them. That is \ref{prop:laplace-uniqueness-sigma-finite}
below. The finite case remains the one almost every caller uses, so it keeps this node and its
tag.}
\end{proposition}

\begin{proof}
\leanok
Substitute $x = e^{-t}$: the map is a measurable embedding carrying a measure on $[0,\infty)$
to one on $(0,1]$, and it turns $\hat\mu(n)$ into the $n$th moment of the image, for every
$n \in \NN$. Two finite measures on the compact $[0,1]$ with the same moments agree, by the
Weierstrass approximation theorem. Transporting back along the embedding loses nothing, so
$\mu = \rho$.
\end{proof}

% additive: not in the draft. The unbounded-measure refinement of A6, split from
% prop:laplace-uniqueness-causal because its hypothesis is genuinely different -- not finiteness
% of the measure but convergence of the transform. Numbered 2.13.
\begin{proposition}[uniqueness without a finiteness assumption]\label{prop:laplace-uniqueness-sigma-finite}
\uses{prop:laplace-uniqueness-causal}
\lean{Hemigroup.laplaceL_injective_of_ne_top}\leanok
Let $\mu, \rho$ be measures on $\RR$ carried by $[0,\infty)$, not assumed finite, and let
$s_0 \in \RR$ be such that $\hat\mu(s_0) < \infty$. If $\hat\mu(s) = \hat\rho(s)$ for every
$s \ge s_0$, then $\mu = \rho$.
\statusT\quad\emph{(\textbf{Proved}, as \texttt{Hemigroup.laplaceL\_injective\_of\_ne\_top},
reducing to Lean core --- so \ledger{A6} stays off the trust boundary in its general form as
well as its restricted one. The hypothesis is worth reading carefully, because it is weaker than
the one the chapter that needs this was expected to supply: what is assumed is not that the
measures are locally finite but that \emph{the transform converges at a single point}. Local
finiteness is then a consequence, \ref{lem:laplace-local-finiteness}, not an assumption.}

\emph{The proof is the classical damping trick and nothing more. $e^{-s_0 t}\mu(dt)$ is finite
precisely because the transform converges at $s_0$; its transform at $s$ is $\hat\mu(s_0 + s)$,
so \ref{prop:laplace-uniqueness-causal} applies to the damped pair; and the damping is undone by
multiplying the density back, which is legitimate because $e^{-s_0t}$ is everywhere positive and
finite. Nothing here needs the tail-agreement form of \ledger{A6} that
\ref{prop:laplace-uniqueness-causal} declines to claim.)}
\end{proposition}

\begin{proof}
\leanok
The measure $\mu_0(dt) := e^{-s_0 t}\,\mu(dt)$ is finite, its mass being $\hat\mu(s_0) < \infty$,
and is carried by $[0,\infty)$, being absolutely continuous with respect to $\mu$; likewise
$\rho_0$, whose mass is $\hat\rho(s_0) = \hat\mu(s_0)$. For $s \ge 0$,
$\hat{\mu_0}(s) = \hat\mu(s_0 + s) = \hat\rho(s_0 + s) = \hat{\rho_0}(s)$, so $\mu_0 = \rho_0$ by
\ref{prop:laplace-uniqueness-causal}. Since $e^{-s_0t}\,e^{s_0t} = 1$ pointwise, multiplying both
by the density $e^{s_0 t}$ recovers $\mu = \rho$.
\end{proof}

\begin{lemma}[convergence forces local finiteness]\label{lem:laplace-local-finiteness}
\uses{prop:laplace-uniqueness-sigma-finite}
\lean{Hemigroup.measure_Icc_ne_top_of_laplaceL_ne_top}\leanok
Let $\mu$ be carried by $[0,\infty)$ with $\hat\mu(s_0) < \infty$ for some $s_0 > 0$. Then
$\mu([0,T]) < \infty$ for every $T$.
\statusT\quad\emph{(\textbf{Proved}, as
\texttt{Hemigroup.measure\_Icc\_ne\_top\_of\_laplaceL\_ne\_top}. Stated because it is what makes
the hypothesis of \ref{prop:laplace-uniqueness-sigma-finite} the natural one: ``locally finite''
would have been an extra assumption, and it is instead a consequence.)}
\end{lemma}

\begin{proof}
\leanok
On $[0,T]$ the integrand satisfies $1 \le e^{s_0 T} e^{-s_0 t}$, so
$\mu([0,T]) \le e^{s_0T}\hat\mu(s_0) < \infty$.
\end{proof}

% additive: not in the draft. The [T] refinement of prop:laplace-continuity.
\begin{proposition}[continuity theorem, causal probability case]\label{prop:laplace-continuity-causal}
\uses{prop:laplace-continuity,prop:laplace-uniqueness-causal}
\lean{Hemigroup.tendsto_integral_of_tendsto_laplace}\leanok
Let $\mu_n, \mu$ be probability measures on $\RR$ carried by $[0,\infty)$, and suppose
\begin{enumerate}
\item $\hat\mu_n(s) \to \hat\mu(s)$ for every $s \ge 0$; and
\item the family is \emph{tight}: for every $\eta > 0$ there is a $T$ with
$\mu_n\bigl((T,\infty)\bigr) \le \eta$ for all $n$, and $\mu\bigl((T,\infty)\bigr) \le \eta$.
\end{enumerate}
Then $\mu_n \to \mu$ weakly.
\statusT\quad\emph{(the direction and the case of \ref{prop:laplace-continuity} that
\ref{thm:main-construction} uses --- transforms to measures, probability form --- and it is
\emph{proved}: \texttt{Hemigroup.tendsto\_integral\_of\_tendsto\_laplace},
\texttt{\#print axioms} reducing to Lean core, so \ledger{A5} does not appear in
\texttt{blueprint/trust-boundary.txt} either. Tightness is a \emph{hypothesis} here rather than a
conclusion, which is the honest shape: \ref{prop:laplace-continuity}'s assignment clause always
held the tightness argument \textbf{[T]}, and \ref{lem:transform-tightness} is where it is
discharged. Feller's boundedness hypothesis, the classical trap of the general form, is carried
for free by the probability form and so cannot be forgotten. The general measure form is
deliberately not formalised.)}
\end{proposition}

\begin{proof}
\leanok
Push forward along $x = e^{-t}$ onto the compact $[0,1]$, as in
\ref{prop:laplace-uniqueness-causal}. There, convergence of the Laplace transforms at the
naturals \emph{is} convergence of the moments, so the images converge weakly by Weierstrass; no
compactness argument is needed, because the carrier is already compact.

Transporting back is where tightness is spent. A bounded continuous $f$ on $[0,\infty)$ pulls
back to $f \circ (-\log)$, which is unbounded at the origin, so it is replaced by
$f \circ \bigl(\min(-\log, T)\bigr)$, bounded and continuous on $[0,1]$ and agreeing with it off
$(T,\infty)$; the two integrals differ by at most $2\|f\|_\infty\,\mu_n\bigl((T,\infty)\bigr)$,
which hypothesis (2) makes uniformly small. An $\varepsilon/3$ argument concludes.
\end{proof}

% additive: not in the draft. The tightness half prop:laplace-continuity's assignment clause
% always held [T] but never stated. Used by thm:main-construction.
\begin{lemma}[tail bound from the transform]\label{lem:transform-tightness}
\lean{Hemigroup.measureReal_Ioi_mul_le}\leanok
Let $\mu$ be a probability measure on $\RR$ carried by $[0,\infty)$. Then for all $s \ge 0$ and
all $T$,
\[
  \mu\bigl((T,\infty)\bigr)\,\bigl(1 - e^{-sT}\bigr) \;\le\; 1 - \hat\mu(s).
\]
Consequently a family of such measures whose transforms satisfy
$\sup_n \bigl(1 - \hat\mu_n(s)\bigr) \to 0$ as $s \to \zp$ is tight.
\statusT\quad\emph{(this is the tightness argument \ref{prop:laplace-continuity} always claimed
as ours, now stated rather than left inside a proof. \texttt{Hemigroup.measureReal\_Ioi\_mul\_le},
\texttt{\#print axioms} reducing to Lean core; the divided form is
\texttt{measureReal\_Ioi\_le\_div}. No hypothesis on $T$ is needed --- causality supplies the
sign that $T \ge 0$ would have. Its use in \ref{thm:main-construction} is
\texttt{kernel\_tail\_le}, which makes the bound uniform over the family as
$\mu_{a,b}\bigl((T,\infty)\bigr) \le F(Bs)/(1 - e^{-sT})$ for $b \le B$, and what makes that
bite is $F(\zp) = 0$; \texttt{isTightMeasureSet\_kernel} repackages it in Mathlib's own
vocabulary, which nothing here now consumes but any other consumer would.)}
\end{lemma}

\begin{proof}
\leanok
Pointwise: $1 - \hat\mu(s) = \int \bigl(1 - e^{-st}\bigr)\,\mu(dt)$, the integrand is at least
$1 - e^{-sT}$ on $(T,\infty)$, and it is nonnegative elsewhere because $\mu$ is carried by
$[0,\infty)$ and $s \ge 0$. Dropping the complement gives the bound.

For the consequence, fix $\eta > 0$. Choose $s > 0$ with
$\sup_n \bigl(1 - \hat\mu_n(s)\bigr) \le \eta/2$, then $T$ large enough that
$1 - e^{-sT} \ge 1/2$; the displayed bound gives $\mu_n\bigl((T,\infty)\bigr) \le \eta$ for
every $n$.
\end{proof}
